\documentclass[11pt]{article}

\usepackage[a4paper,margin=1in]{geometry}
\usepackage{amsmath,amssymb,amsthm,mathtools}
\usepackage{xcolor}
\usepackage{listings}
\usepackage{hyperref}
\usepackage{enumitem}
\usepackage{mdframed}
\usepackage{booktabs}
\usepackage{tabularx}
\usepackage{titlesec}
\usepackage{parskip}

\hypersetup{colorlinks=true, linkcolor=blue!50!black, urlcolor=blue!50!black}

\titlespacing*{\section}{0pt}{16pt}{6pt}
\titlespacing*{\subsection}{0pt}{12pt}{5pt}
\titlespacing*{\subsubsection}{0pt}{9pt}{3pt}
\setlist{nosep, topsep=3pt, itemsep=3pt}
\setlength{\parskip}{6pt}

\definecolor{codebg}{RGB}{243,242,239}
\definecolor{codekeyword}{RGB}{42,120,214}
\definecolor{codecomment}{RGB}{110,110,110}
\lstset{
  basicstyle=\ttfamily\small, backgroundcolor=\color{codebg},
  keywordstyle=\color{codekeyword}\bfseries, commentstyle=\color{codecomment}\itshape,
  frame=single, breaklines=true, showstringspaces=false, tabsize=4,
  columns=flexible, aboveskip=6pt, belowskip=6pt
}

\newmdenv[backgroundcolor=blue!5, linecolor=blue!40!black, linewidth=0.7pt, roundcorner=3pt,
  innermargin=8pt, outermargin=8pt, innertopmargin=6pt, innerbottommargin=6pt, skipabove=8pt, skipbelow=8pt]{cbox}
\newmdenv[backgroundcolor=green!6, linecolor=green!35!black, linewidth=0.7pt, roundcorner=3pt,
  innermargin=8pt, outermargin=8pt, innertopmargin=6pt, innerbottommargin=6pt, skipabove=8pt, skipbelow=8pt]{keyidea}
\newmdenv[backgroundcolor=gray!7, linecolor=gray!50!black, linewidth=0.6pt, roundcorner=2pt,
  innermargin=8pt, outermargin=8pt, innertopmargin=6pt, innerbottommargin=6pt, skipabove=8pt, skipbelow=8pt]{recap}
\newmdenv[backgroundcolor=orange!8, linecolor=orange!50!black, linewidth=0.6pt, roundcorner=2pt,
  innermargin=8pt, outermargin=8pt, innertopmargin=6pt, innerbottommargin=6pt, skipabove=8pt, skipbelow=8pt]{stepbox}

\title{\textbf{Why the Sum-of-Exponentials Trick Works}\\[0.4em]
\large A Proposal: Replacing Fractional-Derivative Memory with a Fixed Number of Exponential States}
\author{}
\date{\today}

\begin{document}
\maketitle

\begin{cbox}
\textbf{What this document is.} A self-contained explanation of the single mathematical idea that makes \texttt{fracmem} possible: an exact identity that rewrites a slowly-decaying power law as a continuum of exponentials, discretized into a short, fixed-size sum. No prior background in fractional calculus is assumed --- every symbol is defined where it first appears, and every abstract step is followed by real numbers, either hand-computed or produced by actually running the code (\texttt{fracmem\_pypi/src/fracmem/soe.py}). The goal is to make the case that this technique is worth adopting more broadly, not just describe it.
\end{cbox}

\tableofcontents

%=================================================================
\section{The problem this solves}

A fractional-order derivative of order $\alpha$ (a non-integer, e.g.\ $\alpha=0.5$, a ``half derivative'') is, in discrete time, a weighted sum over the \emph{entire past} of a signal:
\[
D^\alpha x(t_k) \;\approx\; h^{-\alpha}\sum_{j=0}^{k} w_j\, x_{k-j}, \qquad
w_0 = 1,\quad w_j = w_{j-1}\Big(1-\frac{1+\alpha}{j}\Big).
\]
Here $k$ is which sample we want the answer at, $j$ is how many steps back we're looking, $x_{k-j}$ is the sample $j$ steps before $k$, $h$ is the time between samples, and $w_j$ is a fixed weight (depends only on $\alpha$ and $j$, never on the data) for how much that sample $j$ steps back still matters. The values $w_j$ are called the \textbf{Gr\"unwald--Letnikov (GL) weights}.

The weights decay, but only as a \textbf{power law}:
\[
w_j \;\sim\; j^{-(1+\alpha)}, \qquad j\to\infty.
\]
A power law decays far too slowly to ever safely truncate the sum. Concretely: an ordinary exponential decay $\rho^j$ (say $\rho=0.9$) has essentially vanished by $j\approx 100$ ($0.9^{100}\approx 2.7\times10^{-5}$); a power law $j^{-1.5}$ at $j=100$ is still $10^{-3}$, and at $j=10{,}000$ is still $10^{-6}$ --- a thousand times smaller, not gone. Every past sample keeps contributing, forever. That means computing $D^\alpha x(t_k)$ exactly costs more with every new sample ($O(k)$ work, $O(k)$ memory for that one output), and the total cost of a length-$n$ signal is $O(n^2)$.

\begin{keyidea}
\textbf{The one-sentence goal of this whole document:} find a way to replace the power-law weight sequence $w_j$ with something that can be tracked using a \emph{fixed} number, updated with a \emph{fixed} amount of work per new sample --- no matter how long the signal has been running.
\end{keyidea}

%=================================================================
\section{The key idea: exponentials remember cheaply, power laws don't}

Consider a single exponential sequence $\lambda^j$ for some fixed $0<\lambda<1$ (e.g.\ $\lambda=0.9$). It has a property no power law has: it can be updated with \emph{one multiply}, using only the previous value:
\[
\lambda^{j+1} \;=\; \lambda \cdot \lambda^{j}.
\]
Check it with real numbers: $\lambda=0.9$, $j=3$: $\lambda^3 = 0.729$, and $\lambda^4 = 0.9\times 0.729 = 0.6561$, which is exactly $0.9^4$. You never need to remember $j$ itself, or recompute from scratch --- just keep one running number and multiply by $\lambda$ each step. This is called the \textbf{semigroup property}, and it is the entire reason exponentials are cheap: the whole infinite history compresses into one running state.

A power law has no such shortcut. There is no fixed number $\rho$ with $j^{-\beta} = \rho \cdot (j-1)^{-\beta}$ for every $j$ (the ratio $\big(\frac{j-1}{j}\big)^\beta$ keeps changing as $j$ grows) --- so a power-law memory genuinely cannot be compressed into a single running state the way an exponential can.

\begin{recap}
\textbf{So the plan is:} write the power-law kernel $w_j$ as a combination of a handful of exponentials, $w_j \approx \sum_{i=1}^p c_i\,\lambda_i^{\,j}$. If we can do that honestly, each exponential term becomes one cheap running state, and the whole long-memory sum collapses to $p$ multiply-adds per sample, forever.
\end{recap}

%=================================================================
\section{The exact identity that makes this possible}

\subsection{The Gamma function, briefly}
The Gamma function generalizes the factorial: $\Gamma(n) = (n-1)!$ for a positive integer $n$ (e.g.\ $\Gamma(1)=0!=1$, $\Gamma(4)=3!=6$), and for any real $\beta>0$ it is defined by the integral
\[
\Gamma(\beta) \;=\; \int_0^\infty x^{\beta-1}e^{-x}\,dx.
\]
That's the only property of $\Gamma$ this document needs.

\subsection{The identity, derived from scratch}
\textbf{Claim:} for any $\beta>0$ and any $j>0$,
\[
j^{-\beta} \;=\; \frac{1}{\Gamma(\beta)}\int_0^\infty s^{\beta-1}e^{-sj}\,ds.
\]
\textbf{Derivation.} Start from the definition of $\Gamma(\beta)$ above, and substitute $x = sj$ (so $s = x/j$, $ds = dx/j$):
\[
\Gamma(\beta) = \int_0^\infty x^{\beta-1}e^{-x}\,dx
= \int_0^\infty (sj)^{\beta-1} e^{-sj}\, j\,ds
= j^{\beta}\int_0^\infty s^{\beta-1}e^{-sj}\,ds.
\]
Divide both sides by $j^\beta$:
\[
\frac{\Gamma(\beta)}{j^\beta} = \int_0^\infty s^{\beta-1}e^{-sj}\,ds
\qquad\Longrightarrow\qquad
j^{-\beta} = \frac{1}{\Gamma(\beta)}\int_0^\infty s^{\beta-1}e^{-sj}\,ds. \qquad\blacksquare
\]

\begin{keyidea}
\textbf{Read this identity in words:} the power law $j^{-\beta}$ (hard to track recursively) is \emph{exactly} equal to a continuous, weighted blend of exponentials $e^{-sj}$ (each one trivially trackable), where $s$ ranges over every possible decay rate from $0$ to $\infty$, weighted by $s^{\beta-1}/\Gamma(\beta)$. Nothing is approximated yet --- this is an exact equality.
\end{keyidea}

For the GL tail, the relevant exponent is $\beta = 1+\alpha$ (matching $w_j\sim j^{-(1+\alpha)}$ above), and this is exactly the identity used in \texttt{fracmem\_pypi/src/fracmem/soe.py}:
\begin{lstlisting}[language=Python]
beta = alpha + 1.0
\end{lstlisting}

%=================================================================
\section{From an infinite continuum to a short, fixed sum}

The identity above still has an infinite continuum of exponentials inside an integral --- not yet a finite, computable sum. Two practical steps turn it into one.

\subsection{Step 1: only integrate over the range of $s$ that matters}
As $s\to 0$, $e^{-sj}$ decays negligibly slowly (an almost-flat, near-permanent memory); as $s\to\infty$, $e^{-sj}$ vanishes almost immediately. Neither extreme contributes anything useful past a point, so the code truncates:
\begin{lstlisting}[language=Python]
s_min = 1.0 / max(m_max, 1)   # m_max: longest lag we still care about
s_max = 8.0                    # e^-8 = 0.000335: negligible after one sample
\end{lstlisting}
$m_{\max}$ is the longest distance-into-the-past the design needs to stay accurate for (e.g.\ 2{,}000 samples). $s_{\min}=1/m_{\max}$ is chosen so the slowest exponential included, $e^{-s_{\min}\cdot m_{\max}} = e^{-1}\approx0.37$, has \emph{just barely} not decayed away by the time you reach the longest lag you asked for --- any slower and it would be indistinguishable from ``never decays,'' which is not representable by a finite sum anyway.

\subsection{Step 2: pick a finite number of representative $s$ values, log-spaced}
With the range $[s_{\min}, s_{\max}]$ fixed, approximate the integral by a finite sum evaluated at $p$ representative points $s_1,\dots,s_p$ (this is called \textbf{quadrature}). The code works in $u=\ln s$ instead of $s$ directly:
\begin{lstlisting}[language=Python]
u_min, u_max = np.log(s_min), np.log(s_max)
du = (u_max - u_min) / p
u = u_min + (np.arange(p) + 0.5) * du     # midpoint of each of the p slices
s = np.exp(u)
\end{lstlisting}
\textbf{Why log-spacing, concretely:} with $m_{\max}=2000$, the range is $s\in[0.0005,\,8]$ --- four orders of magnitude. Evenly-spaced points in $s$ itself would put nearly all $p$ points above $s=1$ and leave the slow, long-memory end (which is exactly the hard part of a power law) almost unrepresented. Evenly-spaced points in $u=\ln s$ instead place points evenly across \emph{scales}: one near $s\approx0.001$, one near $s\approx0.01$, one near $s\approx0.1$, one near $s\approx1$, and so on. That matches how the power law itself behaves --- roughly self-similar across scales --- far better than a linear grid could.

The quadrature midpoint rule, $\int f(u)\,du \approx \Delta u\sum_i f(u_i)$, applied to the identity from Section 3 (after substituting $s=e^u$, so $ds = s\,du$) gives
\[
j^{-\beta} \;\approx\; \frac{1}{\Gamma(\beta)}\sum_{i=1}^p \Delta u\; s_i^{\beta}\, e^{-s_i j}.
\]
Matching this to the target form $w_j \approx \sum_i c_i \lambda_i^{\,j}$ identifies
\[
\boxed{\lambda_i = e^{-s_i}} \qquad\text{and}\qquad \boxed{c_i = \frac{A}{\Gamma(\beta)}\,\Delta u\; s_i^{\beta}\, e^{-s_i L}}
\]
where $A = 1/|\Gamma(-\alpha)|$ is a normalization constant tied to the specific GL tail kernel (not the generic identity), and $L$ is the ``local window'' length --- the tail sum is only ever evaluated for lags $j\ge L$, so the weight is centered there. This is exactly the code:
\begin{lstlisting}[language=Python]
lam = np.exp(-s)
c   = (A / gamma(beta)) * du * (s ** beta) * np.exp(-s * L)
\end{lstlisting}

%=================================================================
\section{A tiny example you can trace by hand}

Take $\alpha=0.5$ (so $\beta=1.5$), $L=32$, $m_{\max}=1000$, and just $p=3$ modes --- small enough to trace every number:
\[
A = \frac{1}{|\Gamma(-0.5)|} = \frac{1}{3.544908} = 0.282095,
\qquad \Gamma(1.5) = 0.886227,
\]
\[
s_{\min}=0.001,\quad s_{\max}=8,\quad u_{\min}=\ln(0.001)=-6.907755,\quad u_{\max}=\ln 8=2.079442,
\]
\[
\Delta u = \frac{2.079442 - (-6.907755)}{3} = 2.995732.
\]

\begin{table}[h]
\centering
\small
\begin{tabular}{c c c c c}
\toprule
$i$ & $u_i = u_{\min}+(i+0.5)\Delta u$ & $s_i=e^{u_i}$ & $\lambda_i = e^{-s_i}$ & $c_i$ \\
\midrule
0 & $-5.409889$ & $0.004472$ & $0.995538$ & $2.4716\times10^{-4}$ \\
1 & $-2.414157$ & $0.089443$ & $0.914441$ & $1.4576\times10^{-3}$ \\
2 & $\phantom{-}0.581575$ & $1.788854$ & $0.167152$ & $3.1459\times10^{-25}$ \\
\bottomrule
\end{tabular}
\end{table}

Every one of these numbers comes straight from the four boxed formulas above --- there is nothing else going on. (Mode $i=2$'s weight is vanishingly small here only because $L=32$ pushes $e^{-s_i L}$ to underflow for a fast-decaying mode; this is expected and harmless, not an error.) With just these three exponentials, the approximation is already
\[
w_j \;\approx\; c_0\lambda_0^{\,j} + c_1\lambda_1^{\,j} + c_2\lambda_2^{\,j},
\]
a finite, three-term stand-in for a power law that nominally needs the entire past.

%=================================================================
\section{What the real code produces (verified, not estimated)}

Three modes is enough to see the mechanism; it is not enough for real accuracy. The library's actual default is $p=16$. Running \texttt{soe\_tail\_kernel(alpha=0.5, L=32, p=16, \ldots)} for real (not hand-computed) gives 16 modes ranging from a fast-decaying $\lambda_{15}=0.002708$ down to an extremely slow $\lambda_0 = 0.999324$, and after one additional refinement step (Section 7), the reconstructed tail matches the true GL weights $w_j$ to:

\begin{table}[h]
\centering
\small
\begin{tabular}{rrrr}
\toprule
lag $j$ & exact $w_j$ & SOE approximation & relative error \\
\midrule
32   & $-1.576933\times10^{-3}$ & $-1.777332\times10^{-3}$ & 12.7\% \\
64   & $-5.542212\times10^{-4}$ & $-5.542923\times10^{-4}$ & 0.013\% \\
128  & $-1.953691\times10^{-4}$ & $-1.954030\times10^{-4}$ & 0.017\% \\
256  & $-6.897189\times10^{-5}$ & $-6.895253\times10^{-5}$ & 0.028\% \\
512  & $-2.436736\times10^{-5}$ & $-2.437471\times10^{-5}$ & 0.030\% \\
1000 & $-8.923968\times10^{-6}$ & $-8.922978\times10^{-6}$ & 0.011\% \\
2000 & $-3.154507\times10^{-6}$ & $-3.152163\times10^{-6}$ & 0.074\% \\
4000 & $-1.115182\times10^{-6}$ & $-1.111387\times10^{-6}$ & 0.340\% \\
\bottomrule
\end{tabular}
\end{table}

Worst-case relative error across the entire tail: \textbf{0.34\%}, with 16 numbers standing in for a sum that would otherwise need every one of thousands of past samples. The one weak point is right at $j=L$ (12.7\% error, one lag past the boundary) --- expected, and addressed directly next.

%=================================================================
\section{The one correction the pure math needs}

The quadrature weights $c_i$ from Section 4 are asymptotically correct (accurate for $j\gg L$), but the true GL kernel isn't a perfect power law right at the boundary $j=L$ --- there's a small, known mismatch. \texttt{soe\_refit\_weights} corrects exactly this: it keeps every $\lambda_i$ from the quadrature untouched (they are never re-derived from data), and re-solves for $c_i$ by ordinary linear least squares against the \emph{exact known} GL weights $w_j$ at a few hundred sample lags:
\[
\min_{c}\ \sum_{m} \Big(\underbrace{\textstyle\sum_i c_i \lambda_i^{\,m}}_{\text{SOE approximation}} - \underbrace{w_{L+m}}_{\text{exact, known}}\Big)^2 \Big/ |w_{L+m}|^2.
\]

\begin{keyidea}
\textbf{This is not machine learning, and it is not fit to any real signal.} $w_j$ is a closed-form, exactly-known sequence (Section 1) --- the refit is comparing the approximation against \emph{known ground truth}, the same way you'd round $\pi$ to a nearby fraction and check the error against the true value of $\pi$, not against a dataset. This step is what takes the worst-case error from a few tens of a percent down to the 0.34\% figure in Section 6.
\end{keyidea}

%=================================================================
\section{Why this pays off: from a memory kernel to a running state}

This is the actual payoff, and it resolves a subtlety worth being precise about. Once $(\lambda_i, c_i)$ are fixed, the deployed filter never touches the formulas above again --- it uses the semigroup property from Section 2 directly. Define a running state per mode,
\[
m_i[k] \;=\; \sum_{j\ge 0} \lambda_i^{\,j}\, x_{k-L-j}.
\]
Because $\lambda_i^{\,j+1} = \lambda_i \cdot \lambda_i^{\,j}$, this satisfies the one-line recursion
\[
\boxed{\,m_i[k] = \lambda_i\, m_i[k-1] + x_{k-L}\,}
\]
--- exactly the code in \texttt{fracmem\_pypi/src/fracmem/filter.py}'s \texttt{diagonal\_recurrence}: \texttt{m = lam * m + x\_delayed[k]}. Note $\lambda_i = e^{-s_i}$ (a \emph{per-sample} decay factor, matching this discrete recursion directly) rather than $s_i$ itself (which would be the right quantity only for a \emph{continuous-time} state equation $\dot m_i = -s_i m_i + \dot x$); the code's discrete-time use of \texttt{lam = np.exp(-s)} is the correct choice for exactly this reason.

The tail estimate is then just $\text{tail}[k] = \sum_i c_i\, m_i[k]$: $p$ multiply-adds and $p$ stored numbers, updated once per sample, forever --- replacing the $O(k)$-and-growing direct sum from Section 1 with fixed $O(p)$ work and $O(p)$ memory, independent of how long the signal has been running.

%=================================================================
\section{Why this is worth proposing}

\begin{itemize}
\item \textbf{Complexity.} A signal of length $n$ costs $O(n^2)$ total to differentiate exactly (Section 1); with SOE it costs $O(n\cdot p)$, with $p$ typically 16--32 --- a constant, chosen once, independent of $n$.
\item \textbf{Memory.} Exact computation needs to remember the entire signal history ($O(n)$, growing forever); SOE needs $O(p)$ numbers, fixed for the life of the device.
\item \textbf{Certifiable, not just fast.} Because the decay rates $\lambda_i$ come from an exact identity rather than a data fit, the worst-case error against the true kernel can be computed \emph{before} ever seeing a real signal (Section 6) --- this is a property empirical/black-box approximations don't generally have.
\item \textbf{Applicability.} Any setting with a fractional-order or long-memory power-law kernel benefits the same way: fractional-order PID control, viscoelastic material models, battery impedance models, and more generally any $t^{-\beta}$-type memory kernel, fractional or not.
\end{itemize}

%=================================================================
\section{Where this lives in the code}

\begin{table}[h]
\centering
\small
\begin{tabularx}{\textwidth}{lX}
\toprule
File / function & What it does \\
\midrule
\texttt{soe.py} : \texttt{soe\_nodes\_and\_weights} & Sections 4--5: pure math, builds $(\lambda_i, c_i^{(0)})$ from the Gamma identity, no data. \\
\texttt{soe.py} : \texttt{soe\_refit\_weights} & Section 7: re-solves for $c_i$ against the exact known kernel (still no signal data). \\
\texttt{soe.py} : \texttt{soe\_tail\_kernel} & Runs both of the above in sequence; this is what \texttt{CompressedFractionalFilter} calls internally. \\
\texttt{soe.py} : \texttt{soe\_tail\_error} & Computes the worst-case relative error number reported in Section 6. \\
\texttt{filter.py} : \texttt{diagonal\_recurrence} & Section 8: the deployed $O(p)$-per-sample recursion. \\
\bottomrule
\end{tabularx}
\end{table}

%=================================================================
\section{Further reading}

This is a specific instance of a well-established family of techniques in numerical analysis for accelerating power-law/fractional-memory kernels via sum-of-exponentials approximations, including fast evaluation of the Caputo fractional derivative (Jiang et al., \url{https://arxiv.org/abs/1511.03453}) and more recent exponential-sum approximations for fractional differential equations (\url{https://arxiv.org/abs/2508.20311}).

\begin{recap}
\textbf{Summary, in one paragraph.} A power-law memory kernel cannot be tracked with a fixed amount of state --- but an exact mathematical identity rewrites it as a continuum of exponentials, which \emph{can}. Discretizing that continuum into a handful of representative exponentials (log-spaced, since the kernel spans many time scales), then correcting the resulting weights against the exact known kernel near the boundary, turns an unbounded-memory power-law sum into a fixed-size, fixed-cost recursive filter --- with a computable worst-case accuracy guarantee, before ever touching real data.
\end{recap}

\end{document}
